\documentclass[10pt,letterpaper]{article}

\usepackage[margin=0.68in]{geometry}
\usepackage{array}
\usepackage{booktabs}
\usepackage{longtable}
\usepackage{tabularx}
\usepackage[table]{xcolor}
\usepackage{titlesec}
\usepackage{fancyhdr}
\usepackage{enumitem}
\usepackage{microtype}
\usepackage[hidelinks]{hyperref}

\definecolor{headingblue}{HTML}{1F4E79}
\definecolor{lightblue}{HTML}{D9EAF7}
\definecolor{softgray}{HTML}{F4F6F8}
\definecolor{darkgray}{HTML}{404040}

\pagestyle{fancy}
\fancyhf{}
\lhead{\small EDCI 59100-016}
\chead{\small Week 5 Lecture Notes}
\rhead{\small Summer 2026}
\cfoot{\small \thepage}
\renewcommand{\headrulewidth}{0.4pt}
\setlength{\headheight}{14pt}

\titleformat{\section}{\normalfont\Large\bfseries\color{headingblue}}{\thesection}{0.55em}{}
\titleformat{\subsection}{\normalfont\large\bfseries\color{headingblue}}{\thesubsection}{0.55em}{}
\titleformat{\subsubsection}{\normalfont\normalsize\bfseries\color{darkgray}}{\thesubsubsection}{0.45em}{}
\titlespacing*{\section}{0pt}{9pt}{4pt}
\titlespacing*{\subsection}{0pt}{7pt}{3pt}
\titlespacing*{\subsubsection}{0pt}{5pt}{2pt}

\setlength{\parindent}{0pt}
\setlength{\parskip}{3.8pt}
\setlist[itemize]{leftmargin=1.2em,itemsep=1.5pt,topsep=2pt}
\setlist[enumerate]{leftmargin=1.35em,itemsep=1.5pt,topsep=2pt}
\renewcommand{\arraystretch}{1.15}
\setlength{\tabcolsep}{4pt}
\emergencystretch=3em

\newcolumntype{L}[1]{>{\raggedright\arraybackslash}p{#1}}
\newcolumntype{Y}{>{\raggedright\arraybackslash}X}
\newcommand{\keyterm}[1]{\textbf{\color{headingblue}#1}}

\begin{document}

\begin{center}
{\Large\bfseries Week 5 Comprehensive Lecture Notes}\\[0.25em]
{\normalsize EDCI 59100-016: Literature Review}\\
{\normalsize Drafting the Manuscript, Scholarly Voice, Peer Review, and Revision}\\[0.25em]
{\normalsize Md Kamruzzaman Kamrul \textbar{} Summer 2026}
\end{center}

\tableofcontents
\newpage

\section{Week 5 Map and Source Order}

Week 5 is the transition from planning to manuscript drafting. Week 4 produced a synthesis matrix, theme development, and organizational outline. Week 5 asks the writer to use that work to draft a cohesive literature review manuscript, strengthen academic voice, integrate sources strategically, and use peer feedback for revision.

The saved Week 5 materials should be read in this order:

\begin{enumerate}
\item \textbf{Week 5 overview page.} Introduces the module, learning objectives, weekly tasks, due dates, and Assignment 9 and 10 summaries.
\item \textbf{Where Synthesis Appears in Literature Reviews.} Explains where synthesis is strongest and how synthesis differs across review methodologies.
\item \textbf{Writing from Your Initial Outline.} Shows how to turn Week 4 outline bullets into synthesized paragraphs and linked sections.
\item \textbf{Integrating Sources and Scholarly Voice.} Refines academic voice, citation grouping, source integration, genre awareness, and paragraph-level polish.
\item \textbf{Writing the Full Literature Review Manuscript.} Gives section-by-section drafting guidance for introduction, methods, discussion/implications, limitations, and conclusion across review types.
\item \textbf{Peer Review and Revision.} Explains how to give and use feedback through a reader-focused process.
\item \textbf{Peer Review in Practice.} Discussion activity requiring draft sharing, peer responses, and a revision plan.
\item \textbf{Assignments 9 and 10.} Writing Lab Appointment 2 and Draft Literature Review Paper, both due Sunday, June 21, 2026 at 11:59 PM Eastern.
\item \textbf{External reading.} Rana, Singh, and Kathuria's POWER framework reading on writing effective literature review papers.
\end{enumerate}

The central Week 5 shift is from: \emph{I have a matrix and outline} to \emph{I can write a manuscript section that makes a clear claim, combines sources, interprets patterns, and prepares for feedback}.

\section{Module Overview}

\subsection{Theme}

The Week 5 module is titled \textbf{Drafting the Manuscript}. It focuses on moving from source organization into a substantial literature review draft. The draft should demonstrate synthesis across studies, include methodology-informed structure, and show meaningful progress toward a complete manuscript. Peer feedback and revision are not add-ons; they are central to the week.

\subsection{Learning Objectives}

By the end of Week 5, students should be able to:

\begin{itemize}
\item develop a substantial draft of a literature review manuscript;
\item integrate research across sources rather than summarize individual studies;
\item organize the manuscript with clear headings, subheadings, and a full outline;
\item establish a focused introduction that communicates topic, purpose, and direction;
\item identify patterns, themes, and tensions across studies;
\item incorporate methodology decisions into developing sections;
\item provide and apply peer feedback;
\item revise in response to peer and writing-support feedback.
\end{itemize}

\subsection{Weekly To-Do List}

\begin{longtable}{L{4.8cm} L{8.8cm}}
\toprule
\rowcolor{lightblue}
\textbf{Item} & \textbf{Purpose}\\
\midrule
\endfirsthead
\toprule
\rowcolor{lightblue}
\textbf{Item} & \textbf{Purpose}\\
\midrule
\endhead
\bottomrule
\endfoot
Where Synthesis Appears in Literature Reviews & Locate where synthesis belongs in a manuscript and how it differs by review type.\\
Writing from Your Initial Outline & Convert Week 4 outline points into synthesized paragraphs and linked sections.\\
Integrating Sources and Scholarly Voice & Polish paragraphs through academic voice, citation strategy, sentence flow, and genre awareness.\\
Writing the Full Literature Review Manuscript & Draft the manuscript sections: introduction, methods, body, discussion, limitations, and conclusion.\\
Reading: POWER framework & Understand parameters and decision elements for effective literature review papers.\\
Peer Review and Revision & Prepare to give and use structured peer feedback.\\
Discussion: Peer Review in Practice & Post a draft, review peers' drafts, and identify one revision after feedback.\\
Assignment 9 and Assignment 10 & Complete Writing Lab Appointment 2 and submit the draft literature review paper.\\
\end{longtable}

\section{Where Synthesis Appears in a Literature Review}

\subsection{Main Principle}

Synthesis does not appear equally in every section of a literature review. It is strongest where studies are discussed together and where the writer explains patterns, relationships, differences, or gaps across the literature. A methods section may describe search strategy and selection criteria with minimal synthesis, while body and discussion sections should make synthesis visible.

\subsection{Synthesis by Section}

\begin{longtable}{L{3.3cm} L{4.3cm} L{2.1cm} L{4.0cm}}
\toprule
\rowcolor{lightblue}
\textbf{Section or writing type} & \textbf{What it does} & \textbf{Synthesis level} & \textbf{What writing should look like}\\
\midrule
\endfirsthead
\toprule
\rowcolor{lightblue}
\textbf{Section or writing type} & \textbf{What it does} & \textbf{Synthesis level} & \textbf{What writing should look like}\\
\midrule
\endhead
\bottomrule
\endfoot
Introduction and conclusion & Introduce topic, summarize key ideas, and frame the contribution. & Moderate & Briefly combine research to establish context and overall trends.\\
Main body sections & Present the literature using outline themes, categories, variables, findings, or conceptual sections. & High & Combine multiple studies, compare findings, and explain relationships.\\
Discussion or interpretation & Explain what the mapped patterns mean and why they matter. & Highest & Interpret findings, connect sections, identify gaps, and explain significance.\\
Evidence displays & Present study characteristics, variables, mapped data, or matrix results. & Low & Organize information; do not make tables do all of the interpretive work.\\
Methodology/process description & Explain how the review was conducted. & Minimal & Describe search strategy, screening, selection, and analysis process.\\
\end{longtable}

The practical rule is simple: tables organize evidence, methods explain process, body paragraphs synthesize evidence, and the discussion interprets the significance of the synthesis.

\subsection{Synthesis by Review Type}

Narrative, integrative, and theoretical/conceptual reviews use synthesis to interpret ideas, themes, relationships, concepts, frameworks, or debates. Scoping reviews use synthesis to map categories, trends, distributions, and gaps. Systematic reviews often report structured findings first, then compare and interpret patterns.

For the user's current scoping-review topic, the strongest synthesis should emphasize how research is distributed across VLM architectures, data modalities, construction tasks, evaluation metrics, safety/progress applications, and underdeveloped gaps. It should not claim to determine which intervention is most effective unless the evidence and review design support that claim.

\section{Writing from the Initial Outline}

\subsection{From Outline to Paragraph}

The Week 4 outline becomes the Week 5 draft by turning each outline section into paragraphs that combine studies around a claim. The process is:

\begin{enumerate}
\item Start with a section heading that names the section's main idea.
\item Turn the heading into a topic sentence or claim.
\item Group 2-4 related studies from the matrix or source notes.
\item Combine studies in sentences rather than listing one study at a time.
\item Add comparison, contrast, agreement, variability, or tension.
\item Interpret what the pattern means.
\item End the paragraph or section with a takeaway.
\end{enumerate}

The difference between weak and strong writing is not citation count. Weak writing says, study A found this, study B found this, study C found this. Strong writing says, across these studies, a pattern appears; however, the pattern changes under certain conditions; therefore, the literature suggests a specific interpretation or gap.

\subsection{Basic and Deeper Synthesis}

Basic synthesis combines studies. Deeper synthesis explains what the studies show together. The lesson recommends adding at least one of the following moves:

\begin{itemize}
\item \textbf{Pattern statement:} what consistently appears across studies.
\item \textbf{Tension or contrast:} where findings differ, conflict, or emphasize different dimensions.
\item \textbf{Condition:} when, why, or for whom results change.
\item \textbf{Implication:} what the pattern means for the topic, field, method, or practice.
\end{itemize}

For example, a basic scoping-review sentence might state that VLM studies address safety inspection and progress reporting. A deeper version would explain that safety studies often move toward contextual hazard interpretation, while progress studies expose temporal-reasoning gaps; together, this pattern suggests that construction VLM research is expanding from object recognition toward task-specific decision support, but unevenly across applications.

\subsection{Connecting Paragraphs and Sections}

Strong literature reviews do not restart with every paragraph. Each paragraph should extend, shift, contrast, or refine the previous one. The course suggests using transitions such as:

\begin{itemize}
\item across studies, researchers consistently find;
\item however, findings differ across contexts;
\item in contrast to earlier research;
\item building on previous findings;
\item this extends earlier work by.
\end{itemize}

Before adding a new paragraph, ask what the previous paragraph did and what the next paragraph must add. A paragraph can show a pattern, add a condition, introduce a contrast, explain a limitation, or move from one theme to another. The connection should be explicit, not assumed.

\section{Integrating Sources and Scholarly Voice}

\subsection{Academic Voice}

Week 5 defines academic voice as a formal, clear, evidence-based style of presenting ideas. Academic voice is not complicated vocabulary. It is the writer's decision to make clear claims, use careful language, support those claims with evidence, and foreground patterns rather than conversation or personal opinion.

Academic voice should be:

\begin{itemize}
\item clear and concise;
\item focused on ideas, patterns, and meaning;
\item objective and evidence-based;
\item confident but appropriately cautious;
\item led by the writer's interpretation, with sources supporting the claim.
\end{itemize}

Preferred verbs and stance markers include \emph{suggests}, \emph{indicates}, \emph{highlights}, \emph{demonstrates}, and \emph{points to}. Avoid unsupported certainty such as \emph{proves} unless the evidence actually warrants it.

\subsection{Genre Awareness}

A literature review is a specific academic genre. It has different expectations from a discussion post, research article report, or annotated bibliography. Genre analysis asks how a text works: how ideas are organized, how sources are used, how the writer addresses an audience, how cautious or confident the voice is, and how much detail is needed.

For literature reviews, the genre expectation is to:

\begin{itemize}
\item focus on patterns across studies rather than individual sources;
\item use multiple sources within sentences when they support the same point;
\item develop connections and relationships among ideas;
\item interpret what the research means;
\item keep the writer's argument visible.
\end{itemize}

\subsection{Citation Strategy}

Academic voice and citation work together. The writer makes a claim, supports it with grouped evidence, and then interprets the meaning. Citation strategy should:

\begin{itemize}
\item group studies that support the same idea;
\item avoid repeating separate citations for every sentence when one grouped citation supports the pattern;
\item place citations where they support the claim;
\item avoid author-by-author listing when the goal is synthesis;
\item use comparison and qualification when studies differ.
\end{itemize}

The course's repeated revision principle is: emphasize the idea first, then use citations as support. A choppy paragraph lists authors. A polished paragraph blends sources into a claim, adds nuance, and interprets the pattern.

\subsection{Paragraph Revision Checklist}

\begin{longtable}{L{4.2cm} L{9.3cm}}
\toprule
\rowcolor{lightblue}
\textbf{Revision target} & \textbf{Question to ask}\\
\midrule
\endfirsthead
\toprule
\rowcolor{lightblue}
\textbf{Revision target} & \textbf{Question to ask}\\
\midrule
\endhead
\bottomrule
\endfoot
Sentence structure & Do too many sentences begin with the same phrase or author name?\\
Source integration & Are studies combined around ideas rather than listed separately?\\
Scholarly voice & Does the paragraph sound evidence-based and confident without being inflated?\\
Citation placement & Are citations grouped strategically and not repeated mechanically?\\
Comparison & Does the paragraph show agreement, contrast, condition, or variability?\\
Interpretation & Does the writer explain what the pattern means?\\
Takeaway & Does the paragraph end with a clear contribution to the section's argument?\\
\end{longtable}

\section{Drafting the Full Literature Review Manuscript}

\subsection{Required Manuscript Parts}

The full manuscript draft should include:

\begin{itemize}
\item introduction;
\item methods section aligned with the literature review methodology;
\item body sections developed from synthesized outline sections;
\item discussion or implications;
\item limitations;
\item conclusion.
\end{itemize}

Some results, discussion, implications, or limitation sections may still be developing. The assignment is a draft, not a final paper. A strong draft is developed enough to support meaningful feedback and revision.

\subsection{Introduction}

The introduction tells readers what the topic is, why it matters, and what the review will do. The lesson gives a four-part structure:

\begin{enumerate}
\item \textbf{Topic context:} explain the broad topic.
\item \textbf{What is known:} synthesize what has already been studied.
\item \textbf{Gap or problem:} explain what is missing, unclear, fragmented, inconsistent, or under-mapped.
\item \textbf{Purpose:} state what the review will do and what it contributes.
\end{enumerate}

Different review types shape the introduction differently:

\begin{longtable}{L{3.6cm} L{9.9cm}}
\toprule
\rowcolor{lightblue}
\textbf{Review type} & \textbf{Introduction emphasis}\\
\midrule
\endfirsthead
\toprule
\rowcolor{lightblue}
\textbf{Review type} & \textbf{Introduction emphasis}\\
\midrule
\endhead
\bottomrule
\endfoot
Narrative & Frame the topic broadly, explain key perspectives or debates, and build a coherent interpretive direction.\\
Theoretical / conceptual & Establish why a concept or framework needs clarification, comparison, extension, or refinement.\\
Integrative & Show why diverse empirical, theoretical, or conceptual evidence needs to be brought together.\\
Scoping & Establish a broad and/or emerging topic, show that the scope is unclear, and justify mapping trends and gaps.\\
Systematic & Define a focused problem, summarize known findings, identify inconsistency or missing comparison, and state a clear question.\\
\end{longtable}

For the user's scoping review, the introduction should emphasize rapid growth of VLM and multimodal AI work in construction, the diversity of safety/progress tasks and technical approaches, and the need to map rather than evaluate a single intervention.

\subsection{Methods}

The methods section explains how sources were identified, selected, and analyzed. It should cover:

\begin{itemize}
\item where sources came from: databases, search engines, snowballing, or other search strategies;
\item how inclusion and exclusion decisions were made;
\item how sources were organized, coded, charted, compared, or synthesized.
\end{itemize}

The lesson gives a simple structure: search strategy, selection criteria, and analysis approach. For scoping reviews, the methods section should explain broad search terms, relevance-based inclusion, source types included, grouping/categorization approach, and focus on trends, gaps, and patterns. It can reference PRISMA-style structure if appropriate.

\subsection{Methods Paragraph Builder}

Week 5 also points students back to their working document. The Methods Paragraph Builder activity asks students to add scaffold sheets to the existing workbook and use earlier search logs, screening decisions, source notes, and synthesis planning.

The scaffold should collect:

\begin{itemize}
\item review methodology;
\item search sources or databases;
\item search terms;
\item inclusion and exclusion criteria;
\item source types;
\item reasons for including sources;
\item notes that help describe the process.
\end{itemize}

The paragraph builder then uses blue input cells for review type, topic/focus, databases/search sources, search terms, time span, extra search strategies, criteria, screening process, organization/analysis approach, and the purpose of the methods choice. The generated paragraph is only a starting point; it must be revised to match the actual topic, method, tone, and manuscript structure.

\subsection{Body Sections}

The body is where synthesized outline sections become the substance of the literature review. Body sections should:

\begin{itemize}
\item be organized by theme, category, variable, methodology, chronology, theory, or finding depending on review type;
\item begin with claims rather than source names;
\item integrate multiple sources within paragraphs;
\item compare agreement, difference, contradiction, and condition;
\item include interpretation after evidence;
\item connect across paragraphs and sections.
\end{itemize}

For the user's VLM scoping review, a strong body structure can continue the Week 4 themes: detection to semantic multimodal reasoning; safety compliance and contextual hazard interpretation; activity tracking and temporal reasoning; and dataset, benchmark, and deployment limitations.

\subsection{Discussion and Implications}

The discussion/implications section moves beyond describing the literature to explaining what the literature means. It should:

\begin{itemize}
\item interpret the literature as a whole;
\item explain patterns, relationships, or themes;
\item connect studies that support, contradict, or extend each other;
\item identify implications for practice, research, or theory;
\item name gaps and future directions.
\end{itemize}

It should not summarize individual articles, repeat earlier sections, or be organized by source. The key move is from \emph{Study 1 says; Study 2 says} to \emph{Across studies, the literature shows}.

\begin{longtable}{L{3.6cm} L{9.9cm}}
\toprule
\rowcolor{lightblue}
\textbf{Review type} & \textbf{Discussion emphasis}\\
\midrule
\endfirsthead
\toprule
\rowcolor{lightblue}
\textbf{Review type} & \textbf{Discussion emphasis}\\
\midrule
\endhead
\bottomrule
\endfoot
Narrative & Interpret ideas and perspectives; explain how viewpoints relate or contrast.\\
Theoretical / conceptual & Clarify, connect, refine, or extend concepts and frameworks.\\
Integrative & Synthesize diverse findings to identify patterns and relationships.\\
Scoping & Describe trends, distribution, categories, and gaps rather than effectiveness.\\
Systematic & Evaluate strength, consistency, and implications of evidence.\\
\end{longtable}

\subsection{Limitations}

The limitations section identifies boundaries and constraints of the review. It should not sound like a confession of failure. It demonstrates professional judgment about what the review can and cannot claim.

The lesson's limitation categories are:

\begin{itemize}
\item scope limitations: what was included or excluded;
\item evidence constraints: limitations in available studies;
\item methodological boundaries: what the method allows or does not allow;
\item impact on interpretation: how limitations shape conclusions;
\item future research: what should happen next.
\end{itemize}

For scoping reviews, a typical limitation is that mapping breadth does not evaluate effectiveness or study quality in the same way a systematic review might. Broad inclusion also increases variation in study design, context, and depth, which limits narrow comparison.

\subsection{Conclusion}

The conclusion closes the review by summarizing key insights and reinforcing the review's contribution. It should:

\begin{itemize}
\item summarize the most important patterns;
\item connect back to the purpose or research question;
\item state what the review contributes;
\item explain why the findings matter;
\item point forward briefly.
\end{itemize}

It should not introduce new evidence, new themes, or a second discussion section. Useful openings include: \emph{Overall, the literature shows}; \emph{Taken together, these findings suggest}; \emph{This review contributes by}; and \emph{These insights highlight}.

\section{Methodology-Specific Drafting Guide}

\begin{longtable}{L{3.0cm} L{3.25cm} L{3.55cm} L{3.55cm}}
\toprule
\rowcolor{lightblue}
\textbf{Review type} & \textbf{Introduction} & \textbf{Methods} & \textbf{Discussion / conclusion}\\
\midrule
\endfirsthead
\toprule
\rowcolor{lightblue}
\textbf{Review type} & \textbf{Introduction} & \textbf{Methods} & \textbf{Discussion / conclusion}\\
\midrule
\endhead
\bottomrule
\endfoot
Narrative & Broad topic, key perspectives, debate, coherent narrative. & Targeted searches, influential sources, interpretive synthesis. & Interpret perspectives; explain how viewpoints relate.\\
Theoretical / conceptual & Concept or framework needs clarification, comparison, or extension. & Conceptual relevance, foundational authors, framework comparison. & Clarify concepts or extend frameworks; identify theoretical contribution.\\
Integrative & Diverse evidence needs integration across disciplines, methods, or perspectives. & Multiple evidence types; relevance across designs; comparison and synthesis. & Identify broader patterns and relationships across diverse evidence.\\
Scoping & Emerging or broad field; scope and direction are unclear. & Broad search, relevance-based inclusion, categorization, mapping of trends and gaps. & Describe distributions, concentrations, gaps, and future research directions.\\
Systematic & Focused problem, known findings, inconsistency, clear question. & Structured search, explicit criteria, screening process, extracted outcomes. & Evaluate strength and consistency of evidence; discuss implications and limits.\\
\end{longtable}

\section{Peer Review and Revision}

\subsection{Purpose of Peer Review}

Peer review is a reader-focused process. It is not primarily about fixing grammar, grading a peer, or rewriting someone else's paper. Its purpose is to help a writer understand how the draft is being interpreted and what revisions would strengthen argument, synthesis, and organization.

\subsection{Reviewer Process}

The course recommends a structured process:

\begin{enumerate}
\item Read the entire draft before commenting.
\item Focus first on higher-order concerns: argument clarity, synthesis, organization, and section logic.
\item Provide balanced feedback: at least one strength, one question, and one specific suggestion.
\item Address sentence-level concerns only after meaning and structure.
\end{enumerate}

The key principle is: prioritize meaning before mechanics.

\subsection{Praise, Question, Suggest}

The Week 5 framework for feedback is:

\begin{itemize}
\item \textbf{Praise:} identify what is working.
\item \textbf{Question:} ask about unclear, underdeveloped, or confusing areas.
\item \textbf{Suggest:} offer a specific revision path.
\end{itemize}

Weak feedback says: \emph{This is unclear} or \emph{Looks good}. Strong feedback says: \emph{This paragraph cites several studies but does not explain how they connect; could you synthesize the findings to show the pattern?}

\subsection{Receiving Feedback}

When receiving feedback:

\begin{enumerate}
\item Pause before reacting.
\item Look for patterns across comments.
\item Group similar feedback.
\item Prioritize higher-order concerns.
\item Turn comments into revision tasks.
\item Remember that the writer remains the author and does not need to accept every suggestion.
\end{enumerate}

Revision is not just editing. It may involve re-envisioning the argument, reorganizing sections, improving transitions, and strengthening synthesis.

\section{Peer Review in Practice Discussion}

\subsection{Discussion Requirements}

The discussion is due Thursday, June 18, 2026 at 11:59 PM. Students work in small groups.

Initial post requirements:

\begin{itemize}
\item post the draft paper in Word document format;
\item remind peers which methodology is being used;
\item identify successes and challenges;
\item identify areas where feedback is desired, such as organization;
\item respond to feedback received on Saturday by identifying one key planned revision.
\end{itemize}

Peer response requirements:

\begin{itemize}
\item review two peers' drafts fully before commenting;
\item provide at least one strength;
\item ask at least one question;
\item offer at least two specific suggestions for improvement;
\item optionally share helpful resources;
\item focus on higher-order concerns such as argument, synthesis, and organization.
\end{itemize}

The discussion emphasizes professional tone, specificity, actionable feedback, and questions that prompt deeper thinking.

\subsection{Examples Visible in the Discussion Page}

The saved discussion page shows draft-sharing examples from peers:

\begin{itemize}
\item A scoping review on preservice teacher education for refugee and displacement-affected learning environments, with feedback requested on organization and clarity of synthesis.
\item A scoping review on signal-vehicle cooperative control at signalized intersections, with feedback requested on synthesis, theme clarity, and technical explanation for readers outside transportation engineering.
\item A narrative/thematic review on perceived reasoning in conversational AI virtual agents in VR, with feedback requested on synthesis, transitions, concise academic writing, and the clarity of the argument.
\end{itemize}

These examples reinforce the Week 5 pattern: students are expected to name their methodology, identify the draft's current strengths, state challenges honestly, and ask for targeted feedback.

\section{Assignments}

\subsection{Assignment 9: Writing Lab Appointment 2}

Assignment 9 is due Sunday, June 21, 2026 at 11:59 PM Eastern. Students must attend a second synchronous Purdue Writing Lab appointment, either in person or online. The purpose is to become comfortable sharing writing and using feedback to improve writing.

Required evidence:

\begin{itemize}
\item the Writing Lab coach report;
\item a student reflection;
\item completion during the week before submission.
\end{itemize}

The Week 5 page links to full instructions in an external Google Doc titled \emph{Writing Lab Appointment}. The saved local page includes the assignment summary but not the full Google Doc contents.

\subsection{Assignment 10: Draft Literature Review Paper}

Assignment 10 is due Sunday, June 21, 2026 at 11:59 PM Eastern. It asks for a substantial draft literature review manuscript that moves from planning to writing.

The draft should include:

\begin{itemize}
\item a full manuscript outline with expected sections represented;
\item a clearly developed introduction;
\item a literature review body that synthesizes research rather than summarizes sources;
\item clear organization with headings and subheadings;
\item evidence of connections across studies, including patterns, themes, or tensions;
\item methodology-related sections or content included or in development;
\item optional results, discussion, implications, or limitations if developed.
\end{itemize}

The required submission is a Word document. The file naming instruction is:
\[
\texttt{Draft Literature Review Plan\_SemYR\_Your LastName\_Your FirstName}
\]

The Week 5 page links to a full external Google Doc titled \emph{Draft Literature Review Paper} for complete instructions and grading rubric. The saved local page includes the assignment overview, submission requirement, and link, but not the full rubric text.

\section{POWER Framework Reading}

The Week 5 reading is \emph{Parameters and Decision Elements of Writing Effective Literature Review Papers: Empirical Evidence from Multiple Stakeholders on POWER Framework} by Rana, Singh, and Kathuria. Public metadata identifies it as a 2023 chapter/article in \emph{Review of Management Literature}, volume 2, pages 1-25, with DOI \url{https://doi.org/10.1108/S2754-586520230000002001}.

The POWER framework organizes quality literature review writing around five parameters:

\begin{enumerate}
\item \textbf{Planning:} defining purpose, scope, review type, audience, and contribution before writing.
\item \textbf{Operationalizing:} making the review process transparent through search, selection, coding, analysis, and methodological decisions.
\item \textbf{Writing:} turning evidence into a coherent scholarly manuscript with clear structure and voice.
\item \textbf{Embedding:} presenting the results and contribution that emerge from the review so they are meaningful to the field.
\item \textbf{Reflecting:} checking the review's quality, fit, limitations, and value for readers, reviewers, and editors.
\end{enumerate}

The public abstract reports that the authors refined 256 responses from writers, reviewers, and editors into 37 decision elements. It also notes high priority for embedding and operationalizing, while planning, writing, and reflecting remain important. This aligns closely with Week 5: students are no longer only collecting sources; they are operationalizing their process, writing a coherent manuscript, embedding a contribution, and preparing to revise through feedback.

\section{Applied Notes for the Current VLM Scoping Review}

For the user's literature review on VLMs and multimodal AI for construction safety and progress monitoring, Week 5 translates directly into the following manuscript moves:

\begin{enumerate}
\item \textbf{Introduction:} define construction-site monitoring, VLMs, multimodal AI, safety/progress applications, and spatiotemporal reasoning; then show why the field needs mapping.
\item \textbf{Methods:} describe databases, search strings, inclusion/exclusion criteria, PRISMA-ScR alignment, backward snowballing, and charting variables.
\item \textbf{Body:} organize by themes from Week 4: detection to semantic reasoning; safety compliance and contextual hazards; progress tracking and temporal reasoning; dataset, benchmark, and deployment limits.
\item \textbf{Discussion:} interpret the mapped patterns: the field is moving toward interactive, context-aware, domain-guided VLM systems but remains limited by data, temporal reasoning, benchmarks, prompt reliability, and real-site validation.
\item \textbf{Limitations:} state that a scoping review maps evidence rather than evaluating effectiveness; acknowledge heterogeneity in study design, metrics, and technology maturity.
\item \textbf{Conclusion:} summarize the field's direction and identify future research needs such as construction-specific benchmarks, stronger temporal datasets, and transparent field validation.
\end{enumerate}

\section{Final Week 5 Takeaways}

\begin{enumerate}
\item Week 5 turns the Week 4 outline into manuscript prose.
\item Synthesis is strongest in body and discussion sections, not in methods.
\item Strong paragraphs begin with claims, group studies, compare patterns, and interpret meaning.
\item Scholarly voice means clear, evidence-based, idea-focused writing; it does not mean complicated wording.
\item Citations should be grouped strategically when sources support the same claim.
\item The full draft needs introduction, methods, synthesized body sections, discussion/implications, limitations, and conclusion, even if some sections are still developing.
\item Peer review should focus first on argument, synthesis, and organization.
\item Revision means rethinking structure and meaning, not just editing sentences.
\item Assignment 9 and Assignment 10 are both due Sunday, June 21, 2026 at 11:59 PM Eastern.
\item The POWER framework supports the same logic: plan, operationalize, write, embed contribution, and reflect through revision.
\end{enumerate}

\section*{Sources Used}
\small
Week 5 saved Brightspace pages and embedded lesson pages from the local course folder; Week 5 Peer Review in Practice discussion page; Rana, S., Singh, J., and Kathuria, S. (2023). \emph{Parameters and Decision Elements of Writing Effective Literature Review Papers: Empirical Evidence from Multiple Stakeholders on POWER Framework}. \emph{Review of Management Literature}, 2, 1-25. DOI: \url{https://doi.org/10.1108/S2754-586520230000002001}.

\end{document}
